\section*{Chapter \ref{ch:g6:poverty} --- Poverty and How It Is Counted}

\begin{solution}{exo:g6:poverty:1}
In 1990, a little over $40\,\%$ (the source gives $43.4\,\%$); in 2024, about
$10\,\%$ ($10.1\,\%$).
\end{solution}

\begin{solution}{exo:g6:poverty:2}
Sub-Saharan Africa, at about $45\,\%$ (the source gives $45.0\,\%$).
\end{solution}

\begin{solution}{exo:g6:poverty:3}
$20\,000\,000 \times \frac{15}{100} = 3\,000\,000$: three million people.
\end{solution}

\begin{solution}{exo:g6:poverty:4}
(a) $0.80 + 0.40 + 0.20 + 0.90 + 0.30 + 0.40 = 3.00$ dalens a day.
(b) $3.00 \times 365 = 1\,095$ dalens a year.
\end{solution}

\begin{solution}{exo:g6:poverty:5}
(a) $2.300 - 0.822 = 1.478$ billion people fewer. (b) $1.478 \div 2.300 \approx
0.643$: a fall of $64.3\,\%$.
\end{solution}

\begin{solution}{exo:g6:poverty:6}
(a) Around 1993 (2.34 billion). (b) Yes: the rate had fallen from $47.1\,\%$ in
1981 to $42.1\,\%$ in 1993. The number still rose because the world's
population grew faster than the share fell: a slightly smaller share of a much
larger population can be a larger number.
\end{solution}

\begin{solution}{exo:g6:poverty:7}
(a) Invented example: at the first date, 100 people live on \$2.00 a day; at the
second, 40 of them have moved above the line to \$4.00 a day and 60 still live on
\$2.00. The number of poor fell from 100 to 60, but those still poor are exactly as
poor as before. (b) How far below the line the poor are, date by date: for
example the average income of the poor, or the number living on less than
\$1.50 a day.
\end{solution}

\begin{solution}{exo:g6:poverty:8}
(a) $40 \div 5 = 8.00$ dollars per person per day. (b) Above the \$3.00 line and
above the \$4.20 line, but below the \$8.30 line.
\end{solution}

\begin{solution}{exo:g6:poverty:9}
(a) $100 - 46.2 = 53.8\,\%$. (b) $46.2 - 18.2 = 28.0\,\%$: those below \$8.30 but
not below \$4.20.
\end{solution}

\begin{solution}{exo:g6:poverty:10}
(a) First date: $10 \times 0.30 = 3$ million poor people. Twenty years later:
$20 \times 0.20 = 4$ million. (b) The rate fell (from $30\,\%$ to $20\,\%$), so a
person chosen at random is less likely to be poor: ``it fell''. The number rose
(from 3 to 4 million), so there are more poor people to help: ``it rose''. Both
are true; each answers a different question.
\end{solution}

\begin{solution}{exo:g6:poverty:11}
(a) 1990: $320 \div 2\,300 \approx 0.139$, that is $13.9\,\%$. 2024: $580 \div 822
\approx 0.706$, that is $70.6\,\%$. (b) The number of poor people fell very fast
elsewhere (in East and South Asia above all), while in Sub-Saharan Africa the rate
fell slowly and the population more than doubled, so its number of poor rose. A
region's share of the world's poor depends on the other regions too.
\end{solution}

\begin{solution}{exo:g6:poverty:12}
The value of the line in both years, with how it was set; and the count of the
second year computed again with the first year's line (or the reverse). If, with
the same line, the rate barely moved, the fall came from the change of line. One
can also look at other measures of how people live --- food eaten, income of the
poorest households --- which do not depend on the line.
\end{solution}

\begin{solution}{pb:g6:poverty:1}
\textbf{1.} Bread $200 \times 2.71 = 542.00$; rent $3 \times 1\,000.14 = 3\,000.42$;
transport $50 \times 3.91 = 195.50$; clothing $5 \times 49.00 = 245.00$; energy
$1\,000 \times 0.46 = 460.00$ dalens.
\textbf{2.} $542.00 + 3\,000.42 + 195.50 + 245.00 + 460.00 = 4\,442.92$ dalens a year;
$4\,442.92 \div 365 \approx 12.17$ dalens a day.
\textbf{3.} The basket leaves out needs that are hard to list one by one: water,
soap, medicine, school things; prices also differ from town to town. A round
line above the basket leaves room for them.
\textbf{4.} $5\,000 \div 365 \approx 13.70$ dalens a day.
\textbf{5.} $105 + 126 = 231$ households; $387 + 530 = 917$ people.
\textbf{6.} $387 + 530 + 591 + 967 + 2\,424 = 4\,899$ people. $917 \div 4\,899
\approx 0.187$: a \omterm{def:g6:poverty:poverty-rate}{poverty rate} of $18.7\,\%$.
\textbf{7.} $231 \div 2\,000 = 0.1155$, that is $11.55\,\%$ of households. It is
smaller because poor households are larger than the others: $917 \div 231 \approx 4$
people each, against $4\,899 \div 2\,000 \approx 2.4$ in the whole survey.
\textbf{8.} $11.55\,\%$ is the share of \emph{households}, not of people. The share
of Brenlanders (people) below the line is $18.7\,\%$.
\textbf{9.} $5\,000 \times 1.2 = 6\,000$ dalens a year.
\textbf{10.} $917 + 591 = 1\,508$ people; $1\,508 \div 4\,899 \approx 0.308$, a rate
of $30.8\,\%$.
\textbf{11.} $(1\,508 - 917) \div 917 = 591 \div 917 \approx 0.644$: a rise of
$64.4\,\%$. Many people live just above the old line, between 5\,000 and 6\,000
dalens; a small move of the line counts all of them.
\textbf{12.} For example: how far below the line each poor person is (the table
groups them in bands); how income is shared inside a household; whether the same
households are poor every year; whether they have access to water, schools and
clinics.
\end{solution}
